%% =====================================================================
%%  informe.tex  --  MAESTRO del informe (plantilla del despacho)
%% ---------------------------------------------------------------------
%%  1. Copia esta carpeta:  documentos/informes/YYYY-MM-DD-titulo/
%%  2. Edita metadata.tex y el cuerpo (secciones con \seccion{...}).
%%  3. Compila:  scripts/build.sh <carpeta>/informe.tex
%%
%%  El folio "Pag. x de y" viene activado por defecto en los informes.
%% =====================================================================
\documentclass{informe}

%% =====================================================================
%%  metadata.tex  --  DATOS del informe
%% =====================================================================

\setmeta{numero}{001}
\setmeta{anio}{2026}
% \setmeta{fecha}{04 de agosto del 2026}  % comentado: usa la fecha de hoy

%% El informe lo eleva personal del despacho al titular:
\setmeta{firmante}{Nombre Apellido}
\setmeta{cargofirmante}{Coordinador distrital}


\begin{document}

\numerodocumento

\cabecerainforme[Memorando N°008-2026-2027-\meta{siglas}]
  {\meta{autor} — \meta{cargo}}
  {\meta{firmante} — \meta{cargofirmante}}
  {Resultados de la visita de coordinación al distrito de Quinua}

\seccion{Antecedentes}

Mediante encargo verbal del despacho, con fecha 28 de julio del presente
año, se dispuso la visita de coordinación al distrito de Quinua para el
levantamiento de información sobre el estado de los compromisos asumidos
con las autoridades locales.

\seccion{Análisis}

Durante la visita se sostuvo reuniones con el alcalde distrital y con
los representantes de las comunidades de la zona alta. Los principales
hallazgos son los siguientes:

\begin{itemize}
  \item El expediente técnico del sistema de agua de la zona alta se
        encuentra en revisión desde marzo sin avances.
  \item La institución educativa del centro poblado requiere
        mantenimiento urgente en dos aulas.
  \item Existe interés de la municipalidad en articular una mesa de
        trabajo con el despacho para el mes de setiembre.
\end{itemize}

\seccion{Conclusiones y recomendaciones}

Se recomienda oficiar a la unidad ejecutora correspondiente sobre el
estado del expediente técnico y programar la mesa de trabajo propuesta
dentro de la primera quincena de setiembre.

\atentamente
\bloquefirma[1.8cm]
\adjunto{Panel fotográfico de la visita, 3 hojas.}
\concopia{Archivo del despacho.}

\end{document}
